% nif_groove: a bead in a smooth radial groove on a table turning at constant omega,
% described from two frames (Ch 10, fig:nifgroove).
% Released at rest (relative to the table) at r = a.  Along the groove, in the table's
% frame, x'' = omega^2 x, so x = a cosh(omega t) and v_rel = omega sqrt(x^2 - a^2).
% Seen from the ground the groove has turned through phi = omega t, so the path is the
% spiral r = a cosh(phi), with velocity components (v_r, v_phi) = a omega (sinh phi, cosh phi).
% Drawn with disc radius D = 2.3 cm, a = 0.2 D, bead now at x = 0.5 D (phi = 1.567 rad).
% Forces in the plane of the table, to one scale (1.6 cm per m omega^2 D):
%   centrifugal m omega^2 x = 0.5 m omega^2 D                         -> 0.80 cm
%   wall push N = Coriolis 2 m omega v_rel = 0.917 m omega^2 D        -> 1.47 cm
% The table turns anticlockwise, so N pushes the bead forward (anticlockwise side of
% the groove) and the Coriolis force 2m v_rel x omega points the other way.
% In the ground frame N is the only horizontal force (it supplies a_phi = 2 v_r omega).
% In the table's frame N is balanced across the groove by the Coriolis pseudo-force,
% so the only force ALONG the groove is the centrifugal pseudo-force.
\begin{tikzpicture}[x=1cm,y=1cm]
  \def\D{2.3}
  \def\a{0.2}
  \def\xn{0.5}
  \pgfmathsetmacro\phn{ln(\xn/\a+sqrt((\xn/\a)^2-1))}   % rad, phase now (acosh)
  \pgfmathsetmacro\phe{ln(1/\a+sqrt((1/\a)^2-1))}         % rad, reaches the rim
  \def\gr{15}                                       % groove direction now, left panel
  \pgfmathsetmacro\gz{\gr-deg(\phn)}                % groove direction at release
  \def\kc{0.8}                                     % centrifugal force, cm
  \def\kn{1.466}                                     % N = Coriolis, cm
  % a groove from the centre to the rim along direction #2, disc centred at #1
  \newcommand\groove[2]{%
    \draw[fill=cRock!45,draw=cSteel,thin line,rotate around={#2:#1}]
      ($#1+(0.08,-0.065)$) rectangle ($#1+(\D-0.02,0.065)$);}
  \newcommand\bead[1]{%
    \shade[ball color=cBlock] #1 circle[radius=0.1];
    \draw[cSteel,thin line] #1 circle[radius=0.1];}
  % ============================================================ ground frame
  \coordinate (O1) at (0,0);
  \draw[outline,fill=cSlope] (O1) circle[radius=\D];
  % rotation of the table
  \draw[rotation] (O1)++(70:\D+0.3) arc[start angle=70,end angle=125,radius=\D+0.3];
  \node[anchor=south] at (97:\D+0.35) {$\omega$};
  % where the groove was at release
  \draw[guide] (O1) -- (\gz:\D);
  \groove{(O1)}{\gr}
  % the spiral: travelled (solid), still to come (dashed)
  \draw[cPrim,curve] plot[domain=0:\phn,samples=50,variable=\t]
    ({\gz+deg(\t)}:{\D*\a*cosh(\t)});
  \draw[cPrim,thin line,hidden] plot[domain=\phn:\phe,samples=30,variable=\t]
    ({\gz+deg(\t)}:{\D*\a*cosh(\t)});
  \node[dot] at (O1) {};
  \fill[cPrim] (\gz:\D*\a) circle[radius=0.045];
  \node[cPrim,anchor=north east,inner sep=1.5pt] at (\gz:\D*\a) {start};
  % the bead now: N across the groove, velocity along the spiral
  \coordinate (B1) at (\gr:\xn*\D);
  \pgfmathsetmacro\vang{\gr+atan2(cosh(\phn),sinh(\phn))}
  \draw[velocity] (B1) -- ++(\vang:1.05) node[anchor=west,inner sep=2pt] {$\vv$};
  \draw[force] (B1) -- ++(\gr+90:\kn) node[anchor=south,inner sep=2pt] {$N$};
  \bead{(B1)}
  \node[anchor=north,align=center] at (0,-\D-0.25) {ground frame (inertial):\\ the table turns; the path is a spiral};
  % ============================================================ table's frame
  \coordinate (O2) at (7.4,0);
  \draw[outline,fill=cSlope] (O2) circle[radius=\D];
  \groove{(O2)}{0}
  \node[dot] at (O2) {};
  \coordinate (B2) at ($(O2)+(\xn*\D,0)$);
  % the straight path along the groove
  \draw[cPrim,curve] ($(O2)+(\a*\D,0)$) -- (B2);
  \fill[cPrim] ($(O2)+(\a*\D,0)$) circle[radius=0.045];
  % distance from the axis
  \draw[dim] ($(O2)+(0,-0.3)$) -- node[anchor=north,cSteel] {$x$} ($(B2)+(0,-0.3)$);
  % forces: the real wall push, and the two pseudo-forces (dashed)
  \draw[force] (B2) -- ++(90:\kn) node[anchor=south,inner sep=2pt] {$N$};
  \draw[force,resultant] (B2) -- ++(-90:\kn)
    node[pos=0.72,anchor=west,inner sep=3pt] {$2m\omega v_{\text{rel}}$};
  \draw[force,resultant] (B2) -- ++(0:\kc)
    node[anchor=south west,inner sep=2pt,yshift=2pt] {$m\omega^2x$};
  \bead{(B2)}
  \node[anchor=north,align=center] at (7.4,-\D-0.25) {table's frame (rotating):\\ the table is at rest; the bead slides out along $x$};
\end{tikzpicture}
